% =========================
% Packages, macros, theorem envs (arXiv / amsart version)
% =========================

% 1) Encoding & font enc
\PassOptionsToPackage{utf8}{inputenc}
\usepackage{inputenc}
\PassOptionsToPackage{T1}{fontenc}
\usepackage{fontenc}

% 2) Colors
\PassOptionsToPackage{dvipsnames,svgnames,x11names}{xcolor}
\usepackage{xcolor}

% 3) TikZ for diagrams
\usepackage{tikz}
\usetikzlibrary{calc,arrows.meta,decorations.pathreplacing,patterns,positioning,shapes,fit,intersections}

% 4) Core math / layout packages
\usepackage[margin=1in]{geometry}
\usepackage{mathrsfs}
\usepackage[colorinlistoftodos,textsize=small]{todonotes}

% Inline working-draft markers (remove before submission)
% Use {\color{...} ...} so the macro arguments can span paragraphs/itemize.
\newcommand{\GAP}[1]{{\color{red}\textbf{[GAP]} #1}}
\newcommand{\NEED}[1]{{\color{orange}\textbf{[NEED]} #1}}
\newcommand{\CHECK}[1]{{\color{blue}\textbf{[CHECK]} #1}}
\newcommand{\IDEA}[1]{{\color{teal}\textbf{[IDEA]} #1}}
\PassOptionsToPackage{english}{babel}
\usepackage{babel}
\usepackage{enumitem}
\usepackage{amsmath,amsthm,amssymb}
\usepackage{stmaryrd}
\usepackage{thmtools}
\usepackage{graphicx}
\usepackage{float}
\usepackage{xspace}
\usepackage{calc}

% 5) Hyperref + link colors
\providecolor{Maroon}{cmyk}{0,0.87,0.68,0.32}
\providecolor{RoyalBlue}{cmyk}{1,0.50,0,0}
\usepackage{hyperref}
\hypersetup{
  colorlinks=true,
  breaklinks=true,
  pageanchor=true,
  plainpages=false,
  bookmarksnumbered,
  bookmarksopen=true,
  bookmarksopenlevel=1,
  urlcolor=Maroon,
  linkcolor=RoyalBlue,
  citecolor=Maroon,
  pdfcreator={pdfLaTeX}
}

% 6) Autoref names
\makeatletter
\@ifpackageloaded{babel}{%
  \addto\extrasenglish{%
    \renewcommand*{\figureautorefname}{Figure}%
    \renewcommand*{\tableautorefname}{Table}%
    \renewcommand*{\sectionautorefname}{Section}%
    \renewcommand*{\subsectionautorefname}{Section}%
    \renewcommand*{\subsubsectionautorefname}{Section}%
    \renewcommand*{\theoremautorefname}{Theorem}%
    \renewcommand*{\lemmaautorefname}{Lemma}%
    \renewcommand*{\conjectureautorefname}{Conjecture}%
    \renewcommand*{\corollaryautorefname}{Corollary}%
    \renewcommand*{\definitionautorefname}{Definition}%
    \renewcommand*{\methodautorefname}{Method}%
    \renewcommand*{\factautorefname}{Fact}%
    \renewcommand*{\problemautorefname}{Problem}%
    \renewcommand*{\questionautorefname}{Question}%
    \renewcommand*{\exampleautorefname}{Example}%
    \renewcommand*{\remarkautorefname}{Remark}%
  }%
}{}
\makeatother

% 7) Equation/figure/table numbering by section
\numberwithin{equation}{section}
\numberwithin{figure}{section}
\numberwithin{table}{section}

% 8) Relax line breaking
\emergencystretch=2em

% ── TikZ diagram palette ──
\colorlet{bindA}{red!70!black}       % primary curve
\colorlet{bindB}{blue!70!black}      % secondary curve
\colorlet{fillA}{red!15}
\colorlet{fillB}{blue!15}
\colorlet{fillC}{orange!20}
\colorlet{hlcol}{cyan!70!blue}
\colorlet{hatchcol}{blue!50}

% -------------------------
% Macros & theorem envs
% -------------------------

% closure helper
\newcommand{\closure}[1]{\overline{#1}}

% Xinze's red notes
\newcommand{\xinze}[1]{{\color{red}[\textbf{Xinze:} #1]}}

% Volume command
\newcommand{\Vol}{\mathrm{Vol}}

% Theorem environments
\newtheorem{theorem}{Theorem}[section]
\newtheorem{proposition}[theorem]{Proposition}
\newtheorem{lemma}[theorem]{Lemma}
\newtheorem{corollary}[theorem]{Corollary}
\newtheorem{conjecture}[theorem]{Conjecture}
\newtheorem{claim}[theorem]{Claim}
\newtheorem{subclaim}[theorem]{Sub-claim}

\theoremstyle{definition}
\newtheorem{definition}[theorem]{Definition}
\newtheorem{notation}[theorem]{Notation}
\newtheorem{method}[theorem]{Method}
\newtheorem{fact}[theorem]{Fact}
\newtheorem{problem}[theorem]{Problem}
\newtheorem{question}[theorem]{Question}
\newtheorem{example}[theorem]{Example}
\newtheorem{hypothesis}[theorem]{Hypothesis}

\theoremstyle{remark}
\newtheorem{remark}[theorem]{Remark}

% Compact lists
\newenvironment{itemize*}
  {\begin{itemize}[topsep=-\parskip+\jot,itemsep=-\parskip-\jot]}
  {\end{itemize}}
\newenvironment{enumerate*}
  {\begin{enumerate}[label=(\alph*),topsep=-\parskip+\jot,itemsep=-\parskip-\jot]}
  {\end{enumerate}}
\newenvironment{enumerate**}
  {\begin{enumerate}[label=(\roman*),topsep=-\parskip+\jot,itemsep=-\parskip-\jot]}
  {\end{enumerate}}
\newenvironment{enumerate***}
  {\begin{enumerate}[label=(\alph*'),topsep=-\parskip+\jot,itemsep=-\parskip-\jot]}
  {\end{enumerate}}

% ==== Math macros ====
\newcommand{\cut}{\mathbf{Cut}}
\newcommand{\mres}{\mathbin{\vrule height 1.6ex depth 0pt width
0.13ex\vrule height 0.13ex depth 0pt width 1.3ex}}
\newcommand{\dga}{\dot{\gamma}}
\newcommand{\ddga}{\ddot{\gamma}}
\newcommand{\dddga}{\dddot{\gamma}}
\newcommand{\sss}{\sum_{i=1}^{n}}
\newcommand{\nvec}{\Vec{\mathbf{n}}}
\newcommand{\tvec}{\Vec{\mathbf{t}}}
\newcommand{\bvec}{\Vec{\mathbf{b}}}
\newcommand{\nvecs}{\Vec{\mathbf{n}}_s}
\newcommand{\tvecs}{\Vec{\mathbf{t}}_s}
\newcommand{\bvecs}{\Vec{\mathbf{b}}_s}
\newcommand{\arccur}{\gamma(s)}
\newcommand{\ds}[1]{\frac{d#1}{ds}}
\newcommand{\pd}[2]{\frac{\partial #1}{\partial #2}}
\newcommand{\diff}[2]{\frac{d #1}{d #2}}
\newcommand{\ddiff}[2]{\frac{d^2 #1}{{d #2}^2}}
\newcommand{\sgn}{\mathbf{sgn}}
\newcommand{\emb}{\mathbf{Emb}}
\newcommand{\re}{\mathfrak{Re}}
\newcommand{\im}{\mathfrak{Im}}
\newcommand{\spn}{\mathbf{span}}
\newcommand{\R}{\mathbb{R}}
\newcommand{\RP}{\mathbb{RP}}
\newcommand{\Z}{\mathbb{Z}}
\newcommand{\Q}{\mathbb{Q}}
\newcommand{\N}{\mathbb{N}}
\newcommand{\Om}{\Omega}
\newcommand{\Si}{\Sigma}
\newcommand{\G}{\Gamma}
\newcommand{\haus}{\mathcal{H}}
\DeclareMathOperator{\diam}{\textbf{diam}}
\DeclareMathOperator{\rad}{\textbf{rad}}
\newcommand{\curt}[1]{\textlbrackdbl\#\textrbrackdbl}
\newcommand{\supp}{\mathbf{supp}\,}
\newcommand{\ind}{\mathbf{index}}
\newcommand{\br}[1]{\left\{#1\right\}}
\newcommand{\inp}[1]{\langle #1 \rangle}
\newcommand{\brac}[1]{\left(#1\right)}
\newcommand{\rms}[1]{\mathcal{M}^{k}}
\newcommand{\edge}{\mathscr{E}}
\newcommand{\vertex}{\mathscr{V}}
\newcommand{\eprod}[2]{\langle #1, #2\rangle}
\newcommand{\length}{\mathbf{Length}}
\newcommand{\grad}{\mathbf{grad}}
\newcommand{\id}{\mathbf{Id}}
\newcommand{\dive}{\mathbf{div}}
\newcommand{\bdy}{\partial^*}
\DeclareMathOperator{\per}{per}
\newcommand{\hess}{\mathbf{Hess}}
\newcommand{\Hess}{\mathbf{Hess}}
\newcommand{\ric}{\mathbf{Ric}}
\newcommand{\sect}{\mathbf{Sec}}
\newcommand{\scal}{\mathbf{Scal}}
\newcommand{\tr}{\mathbf{Trace}}
\renewcommand{\det}[1]{\mathbf{det}{(#1})}
\newcommand{\vol}{\mathbf{Vol}}
\newcommand{\dist}{\mathbf{dist}}
\newcommand{\graph}[1]{\mathbf{graph}({#1})}
\newcommand\numberthis{\addtocounter{equation}{1}\tag{\theequation}}
\newcommand{\dat}[2]{\frac{d}{d{#1}}\Bigr|_{{#1} = {#2}}}
\newcommand{\ddat}[2]{\frac{d^2}{d{#1}^2}\Bigr|_{{#1} = {#2}}}
\newcommand{\pik}{\varpi^\kappa}
\newcommand{\Zk}{Z_k(M, \partial M; G)}
\newcommand{\Ztworel}{Z_2(\Omega, \partial \Omega; \mathbb{Z}_2)}
\newcommand{\Ztwo}{Z_2(\Omega; \mathbb{Z}_2)}
\newcommand{\F}{\mathcal{F}}
\newcommand{\eps}{\varepsilon}
\newcommand{\modsp}[1]{\mathbb{M}_{#1}^n}
\newcommand{\mdk}{\mathbf{md}_\kappa}
\newcommand{\snk}{\mathbf{sn}_\kappa}
\newcommand{\csk}{\mathbf{cs}_\kappa}
\newcommand{\cs}{\mathbf{cs}}
\newcommand{\ctk}{\mathbf{ct}_\kappa}
\newcommand{\tnk}{\mathbf{tan}_\kappa}
\newcommand{\an}{An}
\newcommand{\ans}{\mathcal{AN}}
\newcommand{\met}{\mathrm{Met}}
\newcommand{\abs}[1]{\left| #1 \right|}
\newcommand{\norm}[1]{\left\lVert #1 \right\rVert}

% Draft/collaboration comments (REMOVE BEFORE FINAL SUBMISSION)
\newcommand{\qst}[1]{\textcolor{red}{(\textbf{Question: }{#1})}}
\newcommand{\cmt}[1]{\textcolor{blue}{(\textbf{Comments: }{#1})}}
\newcommand{\cmtv}[1]{\textcolor{magenta}{(\textbf{Vitali: }{#1})}}
\newcommand{\corrv}[1]{\textcolor{ForestGreen}{{#1}}}
\newcommand{\detail}[2]{\textcolor{teal}{\textbf{Detail of #1:}#2}}

% Brackets and notation
\newcommand{\llrr}[1]{\llbracket #1 \rrbracket}
\newcommand{\op}{\mathbf{Op}}
\newcommand{\sing}{\mathbf{sing}}
\newcommand{\reg}{\mathbf{reg}}
\newcommand{\C}{\mathbb{C}}
\newcommand{\fo}{\mathscr{F}}
\newcommand{\splm}{\mathcal{L}}
\newcommand{\spdo}{\mathbf{Diff}}
\newcommand{\sh}{\mathfrak{Sh}}
\newcommand{\inj}{\mathbf{Injrad}}
\newcommand{\ghto}{\xrightarrow{\mathbf{G-H}}}
\newcommand{\pghto}{\xrightarrow{\text{pointed $\mathbf{G-H}$}}}
\newcommand{\spa}{\mathbf{Span}}
\newcommand{\radg}{\mathfrak{R}}
\newcommand{\modangle}{\Tilde{\measuredangle}}
\newcommand{\modtriangle}{\Tilde{\triangle}}
\newcommand{\mtr}[3]{[\Tilde{#1}\Tilde{#2}\Tilde{#3}]}
\newcommand{\modcvee}{\Tilde{\curlyvee}}
\newcommand{\mangle}{\measuredangle}
\newcommand{\acts}{\curvearrowright}
\newcommand{\rank}{\mathbf{rank}}
\newcommand{\h}{\mathbb{H}}
\newcommand{\alex}{\mathfrak{Alex}}
\newcommand{\area}{\mathbf{Area}}
\newcommand{\gmh}{\underline{H}}
\newcommand{\sn}{\mathbf{sn}}
\newcommand{\md}{\mathbf{md}}
\newcommand{\modspace}[1]{\mathbb{M}^2\brac{#1}}
\def\dashint{\,\ThisStyle{\ensurestackMath{%
  \stackinset{c}{.2\LMpt}{c}{.5\LMpt}{\SavedStyle-}{\SavedStyle\phantom{\int}}}%
  \setbox0=\hbox{$\SavedStyle\int\,$}\kern-\wd0}\int}
\newcommand{\bigzero}{\mbox{\normalfont\Large\bfseries 0}}
\newcommand{\rvline}{\hspace*{-\arraycolsep}\vline\hspace*{-\arraycolsep}}
\newcommand{\isom}{\textbf{Isom}}
\newcommand{\spt}{\mathbf{spt}}
\newcommand{\lip}[1]{\mathbf{Lip}({#1})}
\newcommand{\CR}{\mathfrak{CompRad}}
\newcommand{\curv}{curv}
\newcommand{\leb}{\mathbf{L}}
\newcommand{\vbar}{\underline{v}}
\newcommand{\aut}{\mathbf{Aut}}
\newcommand{\M}{\mathbb{M}}
\newcommand{\cptsub}{\subset\subset}
\newcommand{\fcy}{\mathcal{Z}}
\newcommand{\fltnorm}{\mathscr{F}}
\newcommand{\Per}{\mathbf{Per}}
\newcommand{\mass}{\mathbf{M}}
\newcommand{\vfd}{\mathscr{V}}
\newcommand{\swp}{\mathscr{S}}
\newcommand{\cridom}{\mathbf{m}}
\newcommand{\criset}{\mathbf{C}}
\newcommand{\inner}{\mathscr{I}}
\renewcommand{\outer}{\mathscr{O}}
\newcommand{\weakstar}{\overset{\ast}{\rightharpoonup}}
\newcommand{\weak}{{\rightharpoonup}}
\DeclareMathOperator*{\esssup}{ess\,sup}
\DeclareMathOperator*{\essinf}{ess\,inf}
\DeclareMathOperator*{\limessinf}{\lim\,ess\,inf}
\DeclareMathOperator*{\limesssup}{\lim\,ess\,sup}
\DeclareMathOperator*{\ess}{ess}
\newcommand{\Snm}{\mathcal{S}_{\mathrm{nm}}}
